\documentclass[11pt]{article}
\usepackage[margin=1in]{geometry}
\usepackage{amsmath, amssymb}

\title{Exercise 3.15: Adding a Constant to All Rewards (Continuing Tasks)}
\author{Sutton \& Barto, Section 3.5}
\date{}

\begin{document}
\maketitle

\section*{Problem}

In the gridworld example, rewards are positive for goals, negative for running
into the edge, and zero otherwise. Are the signs of these rewards important, or
only the intervals between them? Prove, using~(3.8), that adding a constant $c$
to all rewards adds a constant $v_c$ to the values of all states, and thus does
not affect the relative values under any policy. What is $v_c$ in terms of $c$
and $\gamma$?

\section*{Solution}

\subsection*{Proof}

The discounted return (3.8) for a continuing task is:
\[
    G_t = \sum_{k=0}^{\infty} \gamma^k R_{t+k+1}
\]

If we add a constant $c$ to every reward, letting $R'_{t+k+1} = R_{t+k+1} + c$:
\begin{align*}
    G'_t &= \sum_{k=0}^{\infty} \gamma^k R'_{t+k+1}
         = \sum_{k=0}^{\infty} \gamma^k (R_{t+k+1} + c) \\[6pt]
         &= \sum_{k=0}^{\infty} \gamma^k R_{t+k+1}
          + c \sum_{k=0}^{\infty} \gamma^k \\[6pt]
         &= G_t + \frac{c}{1 - \gamma}
\end{align*}

Taking the expectation under any policy $\pi$:
\[
    v'_\pi(s) = \mathbb{E}_\pi[G'_t \mid S_t = s]
              = \mathbb{E}_\pi[G_t \mid S_t = s] + \frac{c}{1 - \gamma}
              = v_\pi(s) + \frac{c}{1 - \gamma}
\]

Therefore:
\[
    \boxed{v_c = \frac{c}{1 - \gamma}}
\]

\subsection*{Interpretation}

Since $v_c$ is the same constant for all states, the difference $v'_\pi(s) - v'_\pi(s')
= v_\pi(s) - v_\pi(s')$ is unchanged. Only the intervals (differences) between
reward values matter, not their absolute signs. Shifting all rewards by $c$ simply
shifts the entire value function by $\frac{c}{1-\gamma}$ without affecting which
states are better or worse.

\end{document}
